\chapter*{How to Use This Booklet}
\addcontentsline{toc}{chapter}{How to Use This Booklet}
\markboth{How to Use This Booklet}{How to Use This Booklet}

This booklet is a self-contained excerpt --- one complete, readable
piece of a much larger work, not a fragment that depends on the rest
to make sense. Everything needed to follow the argument from its first
page to its last is here.

A few conventions carry over from the larger project and are worth
knowing before the first chapter:

\begin{notebox}
Passages set off like this one flag something the main text leans on
without stopping to justify: a term worth pinning down, a distinction
easy to blur, or an aside that matters more than its size suggests.
\end{notebox}

\begin{tipbox}
This style marks a practical suggestion --- a different way to read a
passage, or a pointer toward something worth chasing further once the
chapter is finished.
\end{tipbox}

Footnotes and in-text citations point to the References chapter at the
back, where every work actually cited in this booklet is listed in
full. Sections marked as a work in progress are exactly that: the
argument up to that point stands on its own, but the section itself
may be extended in a later edition.

Beyond that, there is no required order beyond the chapters
themselves, and no assumed background this booklet doesn't supply
along the way.
